\section{Outstanding source requirements and mathematical scope}
This part separates three unfulfilled original-source reading requirements
from the boundaries of the mathematical conclusions. None of the
unavailable originals is used as an unchecked premise in Part I:
the actual proofs identify their opened inputs or give the argument
directly. The report remains marked \texttt{INCOMPLETE} while the
prescribed original readings are outstanding.

\subsection{Resolved definition and specialization questions}
The independent-length toric theorem concerns actual analytic seams.
The additional isotopy question exposed at the third review is now
resolved by Theorem~\ref{A1-toric-nondescent}: raw surgery does not
descend to isotopy classes, even among torically compatible cuts in
one fixed four-pointed surface. The proof compares exact cubic
coefficients in two holomorphic plumbing families. Its phase is
specified by an exact analytic rule, not by a numerically evaluated
angle. Choosing the unique geodesic representative first defines the
separate canonical convention already treated in Part A.

Theorem~\ref{A6-tame-specialization} gives the mixed-characteristic
tame-cover comparison for every existing proper log-regular model
and proves independence of that model without assuming a common
regular refinement. Proposition~\ref{A6-intrinsic-tame} defines the
same tame category intrinsically by all-DVR tests before a model is
available. Proposition~\ref{A6-wild-counterexample} refutes
its unrestricted finite-\'etale strengthening. Existence of such a
model in arbitrary mixed-characteristic dimension is not asserted.
Proposition~\ref{A6-node-character} separately proves the cyclotomic
character times branch-interchange sign and refutes the unqualified
parenthesis at 67/049(b). Proposition~\ref{A6-frobenius-counterexample}
refutes preservation of the same strict Isom-anabelian definition
under positive-characteristic specialization, using two explicit
four-punctured lines over $\mathbf F_{16}$ with characteristic-zero
lifts. Equivalence of tame-cover categories is thereby distinguished
from preservation of reconstruction. These three additional scope
questions are resolved in Part A; no unrestricted model-existence
theorem is inferred.

\subsection{B.3, subject 21: the remaining prescribed original article}
The affine henselian rigidity statement and the completion-complement
comparison have complete checked deductions in
Theorems~\ref{thm:B3-21} and~\ref{thm:B3-fujiwara}, from the opened
Stacks, Haine--Holzschuh--Wolf and Bhatt--Mathew sources. In particular,
the latter original paper's noetherian hypothesis is distinguished
from the subsequent nonnoetherian generalization.

The separately requested direct original reading remains unavailable
for Fujiwara, \emph{Theory of tubular neighborhood in
\'etale topology}, Duke Mathematical Journal 80 (1995), pp.~15--57,
Corollary 6.6.4. This locator has been checked in the opened
later sources, not in the original itself \cite{BFujiwara1995}.

Gabber's original pp.~325--326 are now opened. The definition on p.~325
and complete Theorem1 statement on p.~326 give exactly the
arbitrary-ring, arbitrary-ideal torsion-sheaf theorem used in
Theorem~\ref{thm:B3-21}, whose converse is derived from the original
finite-algebra idempotent criterion. This discharges the required
Gabber1994 derivation. The later original \S5, Corollary1 and the
whole original proof are not claimed read; the complete modern Stacks
proof remains separately identified \cite{BGabber1994}.

The retrieval record includes publisher and historical publisher
endpoints, author and institutional searches, DOI and repository
metadata, public web-archive listings and a public archived landing
page, and whole-journal catalogue records. The exact archived Duke
issue marks its PDF and OCR private; no private file was requested.
The archived article landing page denies full-text access, whereas
the current non-PDF response is an anti-bot page. We do not conflate
these distinct failures. No challenge or subscription control was
bypassed. The smallest remaining gate is direct access to the
specified Fujiwara original statement, followed by a
hypothesis-by-hypothesis comparison with the formulation already
proved here.

\subsection{Scope boundaries which are not claimed as solved problems}
The raw inventory of the twenty-eight thesis subjects and nine
\emph{Pour m\'emoire} headings records the strongest specified proved
formulation used here. Broad historical headings are not silently
converted into blanket theorems. In particular, we do not assert an
unrestricted abelian category of mixed motives, resolution in arbitrary
positive characteristic, every question in SGA 2, or a formal meaning
of the undefined category of ``sch\'emas morcel\'es''. The corresponding
positive results, counterexamples and exact scope are stated in Part B.
This is the explicit ``strongest proved statement'' alternative in
task B.3, rather than a claim that the whole historical programme has
been affirmatively settled.

The positive-characteristic all-prime abelian-fourfold conclusion
explicitly uses Broe's 2026 preprint in
Theorem~\ref{thm:B1b-abelian}. We checked the theorem statement and the
deduction required here, and do not claim an independent verification
of the entire preprint. Ancona's published numerical positivity alone
is not represented as proving that conclusion.

The class requested in B.1(b) is the precisely stated domain of the
componentwise-definiteness argument: algebraic Lefschetz projections
and positivity on the homological primitive algebraic classes.
The criterion in Proposition~\ref{prop:B1b-criterion} proves the
deduction on that entire domain. Its hypotheses are separately
verified for the classes identified in Part B. We do not replace
positivity on that homological space by positivity merely on its
numerical quotient. A geometric classification of all isomorphism
types satisfying the criterion is not asserted or introduced as an
additional requirement of the prompt.

\subsection{The remaining named-source delivery requirement}
The prompt separately instructs us to fetch its named historical works.
This is an outstanding delivery requirement even where an independent
proof or an opened replacement proves the mathematics. In particular,
the source audit has identified outstanding original-text readings for
Pfluger and Lieberman, in addition to the Fujiwara article above.
The requested Raynaud thesis, \emph{Faisceaux amples sur les sch\'emas
en groupes et les espaces homog\`enes}, LNM119 (1970), is distinct from
the Picard article. Its original pp.~14--15 and37--38 are now opened:
the exact property $W$, its base-change compatibility, and the
cohomological-flatness criterion give the complete checked deduction in
Theorem~\ref{thm:B3-raynaud-flatness} and the universal base-change
conclusion in Theorem~\ref{thm:B3-picard}. This named-source requirement
is closed. We do not claim to have read the complete book or the
additional purity results on which its Proposition III.1.4 depends.

The originals of Nakamura (1990), Ihara--Nakamura (1997), and
Gabriel (1932) have now been obtained and the relevant mathematical
passages checked. Their exact hypotheses and their limits are recorded
in the corresponding parts and bibliography. These three source
requirements are closed.

The original Lehto--Virtanen sewing theorem, Alling--Greenleaf's
genus-one classification, and Birman's mapping-class chapter have
subsequently been opened and compared directly with the proofs here.
The exact passages are Chapter II, Sections 6--8, pp.~81--100;
Section 1.9, pp.~60--66; and Section 4.1, pp.~148--161, respectively.
Their source requirements are closed. We retain the distinction between
Birman's closed-surface configuration-space sequence and the puncture
and boundary versions used in Proposition~\ref{A6-birman}.

This list is a live snapshot: the report and bibliography record which
originals have subsequently been obtained, the exact passages actually
read, and the remaining inaccessible works. Original SGA4, $4\frac12$,
5 and6 texts have now been retrieved and selected exact passages read.
Lavrentiev's original 1934 text and Bers's original 1974 announcement
have been opened; the bordered conventions are checked in
\cite[\S\S2(F),3(F),4(E)]{ABers1981}. An announcement of a theorem is
not treated as inspection of an unpublished proof. Koebe's original
1907 covering-radius passage is now read in
\cite[pp.~203--205]{AKoebe1907}. The joint Earle--Marden original
author draft dated24August2011 is now opened and its local construction
derived in Proposition~\ref{A1-Earle-Marden-original}; this closes
A.1(c)'s joint-source requirement. The separate2012 chapter remains
unopened, but the prompt did not specify that chapter as the sole
acceptable joint source. We do not call these
named originals optional or silently replace
them by later accounts. For each mathematical conclusion used in Part I,
the proof instead identifies its actual opened source or supplies the
argument directly.
